This article addresses the task of automated digitization of chemical literature using machine vision and image analysis methods. The relevance of the work is due to the fact that a large part of accumulated chemical knowledge is stored in scanned books, articles, and archival documents that contain not only text but also structural formulas, reaction schemes, tables, and other graphical elements; digitizing such material requires specialized processing methods, since standard text recognition alone cannot extract the meaning of graphical chemical objects.
A chemical document processing pipeline is proposed, including adaptive image preprocessing, page segmentation, extraction of text and graphic regions, recognition of text and chemical structures, and postprocessing with conversion of the results into the machine-readable SMILES format.
Experimental evaluation was performed on images of chemical book pages. For text recognition, the best quality (mean $1-\mathtt{WER}$ score of 0.7693) was achieved by a combined scheme applying adaptive image preprocessing followed by language-model-based text normalization, which improved average OCR quality by about 35\% over the baseline configuration. For chemical structure recognition, the MolScribe, DECIMER, and Qwen3VL models were compared: adaptive preprocessing did not yield a universal improvement --- the best result without preprocessing was obtained by MolScribe, whose quality dropped after preprocessing, DECIMER's result stayed the same, and Qwen3VL showed only a slight gain.
The results show that the effectiveness of image preprocessing depends on the specific model and the type of chemical notation, so it should be tuned for each task rather than applied as a universal processing step.
